Now I have all the reference files. Let me analyze the provided LaTeX against the W1–W4 criteria systematically.

**Analysis of issues found:**

**W1 (禁止語・禁止構文):**
- 「ここではまず, 使い方を体験することを優先する」→ 「ここでは〜を優先する」パターン
- 「式の部品を確認しよう」(2箇所) → 「まず〜を確認しよう」パターン
- 「2つの方向から確かめてみよう」→ メタ宣言
- 「ここでは先に結果だけ見ておこう」→ メタ宣言
- 「考え方がわかったところで, 手元の数字に実際に幅をつけてみよう」→ メタ宣言
- 「この章で身につけたことを整理しよう」→ メタ宣言
- 「この章では, 気づいた上で, それが推定にどう効くかを扱う」→ メタ宣言
- 「次は, これを使って具体的に「幅」を作る方法を見よう」→ メタ宣言

**W2 (アンチパターン):**
- A4 filler: 「判断する材料を手にしたことになる」→ 「ことになる」が情報を足していない

**W4 (日本語の自然さ):**
- 「SDとSE.」→ 短い宣言文を劇的な間として使う英語的パターン
- 「ここまでで, 素朴な方法の弱点が2つ見えた.」→ テキスト進行を参照するメタ的表現
- 「ここまでの議論で」→ 同上

以下、修正を反映したLaTeX全文です:

\section{手元のデータから全体を考える}

  \subsection{問い : 図で見えた手がかりを, この先も言えるか}

    前章で図にしてみると, いくつかの手がかりが見えたはずだ.
    ヒストグラムの山はどちらかに偏っているように見える.
    カテゴリー別の箱ひげ図では, あるカテゴリーの値が他より高そうだ.
    散布図には, 右上がりの傾向がありそうに見える.

    では, この「ありそう」は, どこまで信じてよいのだろうか.

    もし同じやり方でもう一度データを集めたら, 同じ手がかりが見えるだろうか.
    それとも, たまたま今回そう見えただけで, 次に集めたらまるで別の形になるかもしれない.
    手元にあるのは全体のほんの一部だ.
    その一部から, 見えない全体について何かを語ろうとしている.

    選挙の開票速報では, 開票率がまだ1\%の段階で「当確」が打たれることがある.
    全体の1\%しか見ていないのに結果を語れるのは, 一部から全体を推し量るための仕組みがあるからだ.

    前章までは, 手元のデータを要約し, 図にして眺めた.
    やったことは「手元にあるものを整理する」ところまでだ.
    ここから先は, 手元にないもの, つまり全体について語るという, 質的に異なる段階に入る.

    この章では, 図で得た手がかりに\textbf{幅}をつけて,
    「全体でもこのあたりだろう」と語る方法を身につける.
    章の終わりには, 手元の平均や割合に対して
    「この範囲で見ている」と数字つきで報告できるようになる.


  \subsection{素朴な方法 : 手元の平均をそのまま使えばいいか}

    手っ取り早いのは, 手元の標本から計算した平均や割合を, そのまま全体の値だと見なすことだ.
    計算も簡単で, 直感にもかなう.

    しかし, この方法には弱点が2つある.

    \paragraph{少ない観測は, 思った以上にばらつく}

    日本人成人男性の平均身長はおよそ171\,cmで, 標準偏差はおよそ6\,cmとされている.
    街で10人をランダムに選んで身長を測り, 平均を計算したとしよう.
    「170〜172\,cmくらいに収まるだろう」と思うかもしれない.

    実際にはどうか.
    10人の平均がどのくらいばらつくかは, 後で紹介する標準誤差（SE）という量で見積もることができる.
    SE $\approx 6/\sqrt{10} \approx 1.9$\,cm.
    ばらつきの幅をSEの2倍で見ると$\pm 3.8$\,cm, つまり167.2〜174.8\,cmの範囲に広がりうる.
    \textbf{7.6\,cmもの幅}だ.
    10人の平均を「全体の平均」と言い切るには, かなり心もとない.

    では100人ならどうか.
    SE $= 6/\sqrt{100} = 0.6$\,cm.
    幅は$\pm 1.2$\,cm, つまり169.8〜172.2\,cm.
    人数を10倍にすると, 幅はおよそ3分の1になる.

    ネット通販でレビューが3件しかない商品の「平均4.8\,★」と, 300件ある商品の「平均4.3\,★」を見比べたことはないだろうか.
    3件の平均は, たまたま高い評価が続いただけかもしれない.
    300件の平均は, ずっと安定している.
    少ない観測から出た数字をうのみにするのは危うい, という感覚は, 日常の経験とも合うはずだ.

    \paragraph{たくさん集めれば安心か}

    数が少ないとばらつくなら, たくさん集めればよいと思うかもしれない.
    ところが, 集め方に偏りがあると, いくら数を増やしても的外れになる.

    \begin{mybox}{ギャラップ vs.\ Literary Digest : 240万人が大はずれした話}
    1936年のアメリカ大統領選挙で, 雑誌 \textit{Literary Digest} は大規模な郵送調査を行った.
    返送されたのは約240万通.
    これだけの数があれば確実に思えるが, 予測は「ランドンの勝利」.
    結果はルーズベルトの圧勝で, 大はずれだった.

    原因は集め方にある.
    \textit{Literary Digest} は電話帳や自動車登録簿から名簿を作ったが, 1930年代にこれらを持っていたのは富裕層に偏っていた.
    さらに, 郵送に返事をする人としない人の間にも偏りがある.
    結果として, 240万人のデータは母集団を代表していなかった.

    一方, ジョージ・ギャラップはわずか約5万人の調査で正しくルーズベルトの勝利を予測した.
    ギャラップは年齢・性別・地域などの構成比を考慮して回答者を集める工夫をしていた.

    \textbf{$n=240$万でも, 偏りがあれば自信満々に間違える}.
    この事件の後, \textit{Literary Digest} は廃刊に追い込まれた.
    \end{mybox}

    朝だけの来店データで「1日の平均」を語る.
    レビューを書きやすい人だけが残ったデータで「全体の評価」を語る.
    「ちゃんと集める」こと自体が, 実は簡単ではない.

    前章で, 外れ値や偏りに図で気づく方法を見た.
    気づいた偏りが推定にどう効くかが, 次の問題だ.

    素朴な方法の弱点は2つだ.
    少ないデータから出た数字はばらつきが大きいこと.
    集め方に偏りがあると, 数を増やしても解決しないこと.

    では, どうすればよいか.
    必要なのは, \textbf{ばらつきの大きさを見積もって, その分だけ幅を持たせて語る}という発想だ.


  \subsection{核となる考え : 取り直せば毎回ちがう, そのばらつきを数にする}

    コンビニの弁当発注を想像してみよう.
    過去1週間の平均販売数は50個だが, 日によって35個の日もあれば70個の日もある.
    「平均50個」だけを頼りに50個を発注すると, 半分くらいの日は足りなくなるだろう.
    1点の数字だけでは, 判断の材料として不十分だ.
    ばらつきを含めて語る必要がある.

    統計量の推定にも同じ構造がある.

    \paragraph{同じ学校で12通りの「平均」}

    ある学校に生徒が200人いて, 全員の通学時間の平均は30分, 標準偏差は15分だとしよう.
    この中から10人をランダムに選んで平均を計算する, という作業を12の班に頼んだとする.

    12班の結果はたとえばこうなる :
    25, 33, 28, 35, 27, 31, 22, 29, 34, 26, 30, 37（分）.
    最小は22分, 最大は37分.
    \textbf{15分も開きがある}.

    「通学時間の平均は22分」と報告する班と「37分」と報告する班が, 同じ学校から出てくる.
    どちらの班もでたらめをしているわけではない.
    ただ, 選んだ10人が違うだけだ.
    1回の結果を「これが答えだ」と言い切る危うさは, この並びを見れば実感できるだろう.

    手元の推定値は, 同じやり方で取り直したときに生まれる多くの値のうちの\textbf{たった1点}にすぎない.
    このように, 1つの値で母集団の母数を見積もることを\textbf{点推定}と呼ぶ.
    だから, 1点の数字を「確定した事実」として扱うのではなく,
    「ばらつきの中の1点」として捉えることが出発点になる.

    「でも, 実際に標本を何度も取り直すなんてできないのでは」と思うかもしれない.
    たしかに, 現実にはふつう1回しかデータを集められない.
    しかし, 「取り直したらどうなるか」を理論的に見積もることはできる.
    手元のデータのばらつき（$s$）とデータ数（$n$）があれば, それだけで計算できるのだ.

    \paragraph{ばらつきの大きさを測る物差し : 標準誤差}

    推定値がばらつきの中の1点だとわかった.
    では, そのばらつきはどのくらいの大きさなのか.
    大きいのか小さいのかがわからなければ, 推定値をどの程度信用してよいか判断できない.
    欲しいのは, 推定値がどのくらいばらつきそうかの\textbf{物差し}だ.

    第2章で学んだ標準偏差（SD）は, データそのもののばらつきを測る量だった.
    しかし今知りたいのは, 「データのばらつき」ではなく「平均のばらつき」だ.
    同じやり方で標本を取り直したとき, 平均がどのくらい異なる値をとりうるか.
    これを測る量が\textbf{標準誤差}（Standard Error, SE）である.

    平均の場合, SEは次の式で見積もれる.

    \begin{defi}{標準誤差（平均の場合）}{se\_mean}
    標本の標準偏差を$s$, 標本サイズを$n$とする.
    標本平均の標準誤差は
    \[
      \mathrm{SE}(\bar{x}) = \frac{s}{\sqrt{n}}
    \]
    で見積もる.
    \end{defi}

    分子の$s$は, 手元のデータの散らばりの大きさだ.
    データの散らばりが大きいほど, 平均も大きくばらつく.
    これは直感に合う.
    分母の$\sqrt{n}$は, データの件数の平方根だ.
    件数が増えるほどSEは小さくなる.
    つまり, たくさん集めるほど平均は安定する.
    これも直感に合う.

    「$n$そのものではなく$\sqrt{n}$で割るのはなぜだろう」と思うかもしれない.
    具体的な数字で確かめてみよう.
    $n$を4倍に増やすと, $\sqrt{n}$は2倍になり, SEは半分に縮む.
    $n$を100倍にしても, SEは10分の1にしかならない.
    精度の向上は, データ数の増加に対して\textbf{緩やか}だ.
    精度を2倍にしたければ, データは4倍必要になる.

    この関係は実務でそのまま使える.
    工場で製品の重量を抜き取り検査しているとき, 「いま$n=25$で幅がこれだけある.
    $n=100$にすれば幅は半分に縮む」と見通せる.
    全数検査のコストをかけずに, 必要な精度を満たす件数を決められるのだ.

    式は$s/\sqrt{n}$の1本だけだ.
    部品の意味がわかっていれば, 忘れても再構成できる.
    「データの散らばりを, 件数の平方根で割る」.

    \paragraph{SDとSEを混同しない}

    SDとSEは名前が似ていて混同しやすいが, 測っているものがまるで違う.

    \textbf{SD}は, データそのものの散らばりを測る.
    個々のデータが平均からどのくらいまちまちか, を表す.

    \textbf{SE}は, 推定値の散らばりを測る.
    同じやり方で標本を取り直したとき, 平均がどのくらい異なる値になるか, を表す.

    使い分けの目安はこうだ.
    個々のデータの多様さを語りたいならSD.
    平均がどのくらい信用できるかを語りたいならSE.

    データが多様（SDが大きい）であっても, $n$を増やせば推定は安定する（SEは小さくなる）.
    逆に, SDが小さくても$n$が極端に少なければ, SEは大きくなりうる.
    SDとSEは異なる側面を伝えている.
    報告書で「平均 $\pm$ SD」と「平均 $\pm$ SE」を取り違えている例は少なくない.
    見分けの目安として, $n$を増やしてもほぼ変わらないのがSD, $n$に応じて小さくなるのがSEだ.

    \paragraph{この計算が成り立つための条件}

    SEの式$s/\sqrt{n}$は, いくつかの前提のもとで成り立つ.

    第一に, \textbf{無作為に選ぶ}こと.
    母集団のどの要素も等しく選ばれる可能性があること.

    第二に, 各観測が互いに影響し合わないこと（\textbf{ほぼ独立}）.
    友人紹介で似た顧客が続けて現れたり, 同じ人が複数回登場する場合, 実質の情報量は見かけの$n$より少なくなる.

    第三に, \textbf{同じ条件}のもとで集められていること.
    キャンペーン期間と通常期間のデータが混ざると, 同じ集団からの標本とは言いにくい.

    前提が崩れると, 数字は出せても当てにならない.
    ただし, 崩れた場合にも対処の方法はある.
    たとえば「独立でない」データにはクラスター補正, 「条件が異なる」場合には層別推定といった手法が知られている.
    本書では前提が成り立つシンプルな場合を扱うが, 前提が怪しいと感じたときの第一歩は「データの集め方を点検すること」だ.
    この点は9章でもう少し詳しく扱う.

    SEという物差しが手に入った.
    これを使えば, 具体的に「幅」を作れる.


  \subsection{やり方 : 数字に「幅」をつけて言う}

    \paragraph{前提}
    以降の計算は, 前節で述べた前提（無作為, ほぼ独立, 同条件）が成り立ち,
    標本サイズがそこそこ大きい（目安として$n \geq 30$程度）場合を想定している.

    \paragraph{幅のつけ方}

    SEで推定値のばらつきの大きさがわかった.
    では, そのばらつきをどう使えば「全体はこのあたりだろう」と語れるか.

    発想はシンプルだ.
    推定値を中心に, SEの何倍かの幅をとる.

    \begin{defi}{信頼区間（95\%の目安）}{ci95}
    推定値の標準誤差を$\mathrm{SE}$とする. 推定値を中心に
    \[
      \text{推定値} \;\pm\; (\text{およそ}\;2) \times \mathrm{SE}
    \]
    の幅をとったものを, 95\%信頼区間の目安とする.
    \end{defi}

    「およそ2」は, 95\%をカバーする目安の数字だ.
    この数字がなぜうまくいくのかは, 次章で「多く集めると分布の形が安定してくる」という事実を確認する際に明らかになる.

    小さな標本では, 「2」をもう少し大きめにする必要がある.
    たとえば$n=20$前後では2.1程度, $n=10$前後では2.3程度になる.
    $n$がさらに小さいとこの数字はもっと大きくなるが,
    表計算ソフトや統計ソフトが自動で調整してくれる.

    新聞の世論調査で「支持率45\%（$\pm$3ポイント）」という表記を見たことがあるかもしれない.
    あの「$\pm$3ポイント」は, まさに今ここで作っている幅と同じものだ.

    \paragraph{95\%の読み方}

    95\%信頼区間ができたとき,
    「この幅の中に全体の本当の値が95\%の確率で入っている」と読みたくなるかもしれない.
    これは最もよくある誤解だ.

    全体の本当の値（母数）は, 未知ではあるが動かない.
    固定された1つの値だ.
    動くのは, 標本を取り直すたびに変わる推定値のほうであり, したがって区間のほうだ.

    同じ手順で標本を取り, 区間を作ることを100回繰り返したとしよう.
    そのうちおよそ95回の区間は本当の値を含み, 残り5回は外す.
    95\%とは, この\textbf{手順の当たりやすさ}のことだ.
    「今回の区間が当たりかどうか」は, 当たったか外れたかの2択であって,
    95\%の確率で当たりというわけではない.
    報告の際は,
    「この手順で幅を作れば, だいたい100回中95回は当たるように設計されている.
    今回の幅はその1本である」という意味合いで添える.

    \paragraph{なぜ95\%なのか}

    では, 95\%という数字自体に何か特別な根拠があるのだろうか.
    ない.
    95\%は慣習だ.
    90\%にすれば幅は狭くなるが, 100回のうち10回は外す.
    99\%にすれば外す回数は減るが, 幅は広がる.
    「どのくらいの外れを許容するか」は, 分析する側が目的に応じて選ぶことであり,
    理論が一意に決めるものではない.
    95\%が多く使われるのは, 幅の狭さと外しにくさのバランスが実用上とりやすいからだ.

    \paragraph{幅で語ると何がうれしいか}

    新薬の治験で「血圧を平均5\,mmHg下げた」という結果が出たとする.
    もし95\%信頼区間が2〜8\,mmHgなら, 「少なくとも2\,mmHgは下がりそうだ」と見込める.
    医師はこの情報をもとに投薬を判断できる.
    ところが区間が$-1$〜$11$\,mmHgなら, 「効かない可能性もある」と判断は変わる.
    同じ「5下がった」でも, 幅次第で意味がまるで違う.

    「商品の平均重量は170\,g」と言い切るのと, 「167〜173\,gの幅で見ている」と添えるのでは, 意思決定の安全度がまるで違う.
    信頼区間は, 過大な断言を防ぎ, 判断ミスを減らすための道具だ.

    では, 手元の数字に実際に幅をつけてみる.


  \subsection{体験 : 手を動かして3つのケース}

    \paragraph{ケース1 : 平均に幅をつける}

    ある商品の重量を$n=64$個測ったところ, 平均$\bar{x} = 170$\,g, 標準偏差$s = 12$\,gだった.

    SEは$12/\sqrt{64} = 12/8 = 1.5$\,g.
    95\%の目安は
    \[
      170 \pm 2 \times 1.5 \;\Rightarrow\; 167.0\;\text{g〜}\;173.0\;\text{g}.
    \]

    「このやり方を繰り返せば, こういう幅がだいたい当たる. 今回の幅はその1本」と読む.

    この幅が「広い」か「狭い」かは, 文脈で決まる.
    たとえば, この商品の規格が「165〜175\,gに収まること」なら,
    幅はすっぽり規格内に入っており安心材料になる.
    逆に, 規格が「169〜171\,g」と厳しければ, 幅がはみ出しており,
    もう少し精度を上げる（$n$を増やすなどの）対策が必要だ.
    幅そのものに「良い」「悪い」はなく, 何と比べるかで評価が変わる.

    \paragraph{ケース2 : 割合に幅をつける}

    20商品のうち, 星評価が4.0以上のものが$k=14$品だったとする.
    割合は$p = 14/20 = 0.70$.

    割合のSEは
    \[
      \mathrm{SE}(p) \approx \sqrt{\frac{p(1-p)}{n}} = \sqrt{\frac{0.70 \times 0.30}{20}} \approx 0.102.
    \]

    分子の$p(1-p)$は割合のばらつきの源だ.
    $p$が0.5に近いほどばらつきは大きく, 0や1に近いほど小さくなる.
    分母の$n$は件数で, 件数が多いほどSEは小さくなる.
    平均のSEと同じ構造だ.
    「推定値のばらつきを, データから見積もる」という考え方は共通している.

    95\%の目安は
    \[
      0.70 \pm 2 \times 0.102 \;\Rightarrow\; 0.50\;\text{〜}\;0.90.
    \]

    かなり広い.
    $n=20$では, 割合の推定はこの程度の精度にとどまる.

    端に寄った割合（0や1に近い）のときは, この式がやや楽観的になりがちだ.
    表計算ソフトや統計ソフトは, その場合に少し広めの安全な幅を自動で作ってくれる.

    \paragraph{ケース3 : 2つの平均の差に幅をつける}

    カテゴリーAとBで値の平均を比べたいとする.
    2群が互いに独立（AとBで別の個体から集めている）であるとき,
    各群の平均$\bar{x}_A, \bar{x}_B$, 標準偏差$s_A, s_B$, 件数$n_A, n_B$を使う.

    差のSEは
    \[
      \mathrm{SE}(\bar{x}_A - \bar{x}_B) \approx \sqrt{\frac{s_A^2}{n_A} + \frac{s_B^2}{n_B}}.
    \]

    それぞれの群の分散を件数で割ったものを足し合わせ, 平方根をとっている.
    2群が独立であれば, ばらつきの大きさ（分散）は足し算できる.
    件数が多い群やばらつきが小さい群は, 差の推定を安定させる方向に寄与する.

    幅は同じく, 差 $\pm$ およそ$2 \times \mathrm{SE}$.

    前章で箱ひげ図を使って2つのカテゴリーの違いを目で確かめた.
    あのときは「違いがありそう」という印象だった.
    ここでは, その「違い」に具体的な幅をつけて,
    偶然のばらつきの範囲内かどうかを判断できるようになった.
    図で「見えた違い」を, 数字で裏づける段階に進んだのだ.

    \paragraph{幅を狭めたいときの3つの方法}

    区間が広すぎて実用にならないと感じたら, 3つの調整が考えられる.

    $n$\textbf{を増やす}.
    SEは$1/\sqrt{n}$のペースで縮む.
    現在$n=20$で幅が広いなら, $n=80$にすれば幅はおよそ半分になる.
    「あと何件集めれば幅がどのくらい縮むか」を, この関係で見積もれる.

    \textbf{ばらつき}$s$\textbf{を下げる}.
    測定条件を安定させたり, 対象を絞ったり, 層に分けて推定したりすると,
    $s$が小さくなりSEも縮む.

    \textbf{信頼水準を下げる}.
    95\%を90\%にすれば幅は狭くなるが, 外しやすくもなる.
    これはトレードオフであり, 目的に応じて選ぶ.

    実務では, 最初の2つ（件数の増加とばらつきの低減）が本筋であることが多い.

    \paragraph{実務で気をつけること}

    SEは$n$を増やせば縮む.
    区間も狭くなる.
    しかし, 集め方の偏りは$n$をいくら増やしても消えない.
    偏ったデータを大量に集めると, 狭い区間が得られるが,
    その区間の中心が的外れになる.
    \textbf{精密に間違える}という, 最も危険な状態だ.
    区間が狭いからといって安心するのではなく, データの集め方をまず点検する.
    この点は9章で改めて整理する.

    平均とSDは外れ値に引っ張られやすい.
    前章でヒストグラムや箱ひげ図から外れ値に気づく方法を扱ったが,
    ここでの問題は, その外れ値が推定の幅に直接影響することだ.
    きつい外れ値がある場合, 中央値やIQRで要約し,
    幅のつけ方も\textbf{ブートストラップ}などの別手法を使うほうが安全なことがある.
    ここでは名前だけの紹介にとどめるが,
    表計算ソフトや統計ソフトの多くがこの機能を備えている.

    推定の幅は, 期間の切り取り方, 欠損値の扱い, 外れ値の処理によって変わる.
    同じデータでも, 前処理の仕方が異なれば幅は変わりうる.
    「何を, いつ, どう集めて, どう処理したか」を明示しておくことが,
    結果を他者と共有するための最低限の作法だ.


  \subsection{結果と次の一手 : 幅で語れたら, 次は根拠を確かめる}

    標本を取り直せば, 平均も割合も毎回異なる値になる.
    手元の数字は, 多くのありえた値の中のたった1点にすぎない.
    だから, 1点の数字ではなく\textbf{幅をつけて}語る.

    幅の物差しが\textbf{標準誤差}（SE）だ.
    平均の場合は$s/\sqrt{n}$で見積もれる.
    幅の作り方は, 推定値 $\pm$ およそ$2 \times \mathrm{SE}$（95\%の目安）.

    この幅があれば,
    「全体の平均はこの範囲で見ている」と,
    過大な断言を避けて語れる.
    幅が広いならデータを増やすかばらつきを減らすかで精度を上げる見通しも立つ.

    「およそ2」という数字を何度か使った.
    この数字は, なぜ95\%のカバーに対応するのだろうか.
    次章では, 標本平均を繰り返し計算したときに現れる分布の\textbf{形}に注目する.
    多く集めると, その形が特定のパターンに安定してくるという事実を確かめる.
    この事実が, 「およそ2でうまくいく」根拠を与えてくれる.
    そしてその先には, 見えた差が偶然のばらつきだけでは説明しにくいと言えるかどうかを問う段階が待っている.


  \subsection*{作業 : ランニングデータで幅をつけて言う（値は後で挿入）}

    ランニング例のデータについて, 次を計算し,
    幅を含めた一文の所見を書け（小数第2位まで, 単位併記）.

    \begin{enumerate}
      \item \textbf{平均の幅}（95\%目安）:
            $\bar{x}, s, n, \mathrm{SE}$を求め,
            $\bar{x} \pm 2 \times \mathrm{SE}$で幅を作れ.
      \item \textbf{星評価\,$\geq 4.0$\,の割合の幅}（95\%目安）:
            $k, n, p = k/n$を求め,
            $\mathrm{SE} \approx \sqrt{p(1-p)/n}$で幅を作れ.
            端に近いときは, 道具が出す幅を優先してよい.
      \item \textbf{カテゴリーAとBの平均差の幅}（95\%目安）:
            差, $s_A, s_B, n_A, n_B$を求め,
            $\mathrm{SE} \approx \sqrt{s_A^2/n_A + s_B^2/n_B}$で幅を作れ.
    \end{enumerate}

    記録の約束 : 対象と期間, 欠損や外れ値の扱い, 切り取り基準をノートに明示すること.
